\documentclass[11pt]{article}
% Uses the canonical latex-docs semantic style in-repository. The local
% lightweight fallback is for independent compilation of this distribution.
\IfFileExists{technical-user-manual.sty}{\usepackage{technical-user-manual}}{\usepackage{../support/aaoa-standalone}}
\usepackage{mathtools}
\setDocTitle{Chapter 9: Words and Mappings}
\setDocSubtitle{Guided reading, derivations and original exercises}
\setDocVersion{1.0}
\setDocStatus{Educational reference / unapproved working material}
\setDocOwner{Algorithm Engineering}
\setDocAuthor{Original companion material}
\setDocSummary{Coverage of primary book Chapter 9 with independently authored examples, proof obligations and lab transfer.}
\providecommand{\Expect}{\mathbb{E}}
\providecommand{\Prob}{\mathbb{P}}
\providecommand{\Var}{\operatorname{Var}}
\providecommand{\bigoh}{\mathcal{O}}
\begin{document}
\makedoctocpage
\section{Reading map and prerequisite check}
\noindent\textbf{Primary source sections:} 9.1 Separate Chaining; 9.2 Balls-and-Urns Model and Properties of Words; 9.3 Birthday Paradox and Coupon Collector; 9.4 Occupancy Restrictions; 9.5 Occupancy Distributions; 9.6 Open Addressing; 9.7 Mappings; 9.8 Integer Factorization and Mappings (printed pp. 474--532).\par
\noindent\textbf{Prerequisites:} Completed preceding mathematical-method chapters or equivalent demonstrated mastery.\par
\section{Learning outcomes}
\begin{enumerate}
\item Derive expected bucket occupancy under independent uniform hashing.
\item Relate occupancy distributions to chaining and birthday/coupon models.
\item Recognize dependence and probing effects in open addressing.
\end{enumerate}
\section{Guided technical analysis}
\subsection{Uniform occupancy hypothesis}
If $m$ independent keys choose one of $b$ buckets uniformly, the load factor is $\lambda=m/b$. For fixed bucket $j$, $X_j=\sum_{i=1}^m I\{h(k_i)=j\}$ has expectation $m/b$ and variance $m(1/b)(1-1/b)$.
\subsection{Collisions and exposure}
The expected number of colliding key pairs is $\binom m2/b$, by applying linearity to each pair. Real-world low-entropy keys, weak hash functions, and adversarial inputs can violate the independent uniform assumption. State whether a measured collision is an equal hash, equal bucket, or failed probe.
\subsection{Other mapping models}
A random mapping from $n$ elements to $n$ elements is not the same as a random permutation: the former need not be bijective. Do not apply permutation-cycle theorems to arbitrary functional digraphs without adjusting the distribution.
\section{Independent problem set}
\begin{enumerate}
\item Derive $\Expect[\#\text{colliding pairs}]=\binom m2/b$ and check by exhaustive enumeration for $m=3,b=2$.
\item Simulate coupon collection for $b=8$ and compare its sample mean to $8H_8$.
\item Measure separate-chaining lengths at several load factors and document deviations from a Poisson approximation.
\end{enumerate}
\section{Laboratory transfer}
\textbf{Required laboratory:} LAB-CPP02: deterministic bucket-placement trials.\par
\noindent\textbf{Minimum recorded evidence:} repository commit or source hash; test inputs/seeds; actual operation counts; predicted values; one explicit model limitation.\par
\section{Mastery gate and self-review}
Raw bucket counts sum to $m$ and reported load factor equals $m/b$; model limitations explicitly recorded.\par
\begin{enumerate}\item Can I define the random object/input exactly? \item Is my identity exact, asymptotic, or empirical? \item Which base cases and counterexamples did I check? \item What would falsify my performance hypothesis?\end{enumerate}
\section{Suggested notebook structure}
Use five named sections: model and invariants; derivation; deterministic small instances; randomized replication; discrepancy and revised hypothesis. Keep a proof separate from measurement results.

\end{document}
